%\chapter*{Anhang}
%\large{
%Anhang 1\par
%Anhang 2


% setzt die Bildunterschriften auf A1 (A=Anhang)
\renewcommand{\thefigure}{A\arabic{figure}} 
\chapter{Anhang}

\subsection*{Elektronischer Anhang }

\subsection*{Fahrzeugdaten}
\begin{figure}[hb!]
	\caption{Engine Power}
	\includegraphics[width=350pt]{Abbildungen/Engine_power.png}
	\label{fig:enginePower}
\end{figure}

\begin{figure}[hb!]
	\caption{Max Tire Force}
	\includegraphics[width=350pt]{Abbildungen/Engine_power.png}
	\label{fig:maxTireForce}
\end{figure}



 \subsection*{Einbindung Grafik im Anhang}
  \begin{figure}[h!]
  \begin{centering}
  {\includegraphics[width=0.33\textwidth]{Abbildungen/fau.png}}
   \caption{Unterschrift Bild x Die auf die Rotationsfrequenz des Innenzylinders normierten Eigenfrequenzen der gefun-denen Grundmoden der Taylor-Strömung für h(Die azimutale Wellenzahl ist mit m  bezeichnet.)}
  \end{centering}
  \end{figure}

\subsection{Code}

\begin{lstlisting}
[caption={ein paar Zeilen code}\label{lst:vehicleModel},captionpos=t] 
    for i in 0: N-1
    
    @NLexpression(m, x_dd, (VehicleModel.max_long_acc * x[8*i + 7]/10.0))
    
    @NLconstraints(m, begin
    x[(i + 1)*8 + 1] - (x[i * 8 + 1] + x[8*i + 3]*
    dt*cos(x[i*8 + 4] + atan(lr/(lf + lr) * tan(x[i*8 + 8])))) == 0
    x[(i + 1)*8 + 2] - (x[i * 8 + 2] + x[8*i + 3]*
    dt*sin(x[i*8 + 4] + atan(lr/(lf + lr) * tan(x[i*8 + 8])))) == 0
    x[(i + 1)*8 + 3] - (x[i * 8 + 3] + x_dd*dt) == 0
    x[(i + 1)*8 + 4] - (x[i * 8 + 4] + x[8*i + 3]*
    dt / lr*sin(atan(lr/(lf + lr) * tan(x[i*8 + 8])))) == 0
    atan(0.5 * (lf + lr) * mpc_struct.beta_max / x[i*8 + 3]^2)
    - atan(lr/(lf + lf) * tan(x[i*8 + 8])) >= 0  #max_beta - beta
    atan(0.5 * (lf + lr) * mpc_struct.beta_max / x[i*8 + 3]^2)
    + atan(lr/(lf + lf) * tan(x[i*8 + 8])) >= 0  #max_beta + beta
    end)
    end
\end{lstlisting}





